\section{Related Work}\label{sec:related}

\paragraph{Vision-Language-Action Models.}
VLA models combine pretrained vision-language backbones with action prediction for generalist manipulation~\cite{brohan2023rt2,black2024pi0,kim2024openvla}. Recent work enhances spatial reasoning through auxiliary 3D supervision~\cite{li2025spatialforcing,li2025qdepth,depthvla2025}, but these methods target perception rather than safety, with no collision constraints. Pi0~\cite{black2024pi0} and its successor Pi0.5 use flow matching for action generation, where near-zero actions can arise from training instabilities~\cite{dong2025conditioning}. More broadly, diffusion policies have been shown to operate as action lookup tables that memorize training sequences rather than generalizing~\cite{he2025demystifying}, providing a mechanistic explanation for action collapse when the policy encounters out-of-distribution states. We test both families to separate architecture effects from safety evaluation artifacts. Whether geometric supervision can teach active obstacle avoidance is the question our work addresses.

\paragraph{Safety in Robot Learning.}
Safety-aware approaches span training-time constrained RL~\cite{zhang2025safevla} and test-time filters: control barrier functions~\cite{hu2025vlsa,singletary2022safecritical}, collision risk gating~\cite{cofreeevla2026}, inverse kinematics QP~\cite{spalter2024ikinqp}, and Diffusion Policy collision costs~\cite{li2024langguided,yang2025cape}. SafeDiffuser~\cite{xiao2025safediffuser} integrates barrier functions into diffusion denoising but does not learn intrinsic collision awareness. A recent survey~\cite{li2026vlasafety} categorizes these into runtime shields, constrained optimization, and representation-based methods. Common to all is that safety is imposed extrinsically, adding inference-time overhead and requiring collision geometry at deployment. This work tests whether auxiliary SDF supervision can instead produce intrinsic collision awareness.

\paragraph{Signed Distance Functions in Robotics.}
SDFs encode minimum surface distance, with sign indicating free space or penetration, extending the artificial potential field paradigm~\cite{khatib1986potential} to continuous shape representations. Neural SDFs~\cite{park2019deepsdf} enable continuous shape representation; in robotics, SDFs serve perception~\cite{ortiz2022isdf}, planning~\cite{vasilopoulos2023ramp}, and differentiable collision detection~\cite{zhi2021diffco}. CDF~\cite{li2024cdf} and CSSDF-Net~\cite{chen2026cssdfnet} extend SDFs to configuration space for manipulation planning. RDF~\cite{li2024rdf} represents robot geometry as composable per-link SDFs, enabling whole-body collision queries, directly relevant to the full-body gap our work identifies. Neural SDF-based MPC~\cite{jacquet2025neuralnmpc} is effective in classical pipelines, but integration into end-to-end learned policies remains unexplored. All successful SDF applications in robotics operate on full-body or configuration-space representations; restricting supervision to the EE does not transfer these benefits.

\paragraph{SDF integration pathways.} SDF-based geometric supervision can enter VLA policies through three mechanisms: (a)~\emph{training-time auxiliary loss} on the EE trajectory (tested here); (b)~\emph{constrained optimization}, where SafeVLA~\cite{zhang2025safevla} embeds safety constraints via Lagrangian optimization; and (c)~\emph{runtime filtering}, where VLSA~\cite{hu2025vlsa} applies control barrier functions at inference. Pathway~(a) provides gradients only at the supervised point, leaving intermediate links unmonitored; (b) and~(c) can query full-body geometry, sidestepping this bottleneck. Our negative finding about training-time SDF is specific to~(a); our oracle experiments (\S\ref{sec:ee_only}) additionally test a restricted, EE-only form of runtime refinement.

\paragraph{Evaluation Metrics for Safe Manipulation.}
SafeLIBERO~\cite{hu2025vlsa} extends LIBERO~\cite{liu2023libero} with obstacle-augmented environments and a CAR metric. However, CAR treats all collision-free episodes identically: an action-collapsed robot scores near-perfect CAR without any avoidance behavior, echoing concerns about metric reliability~\cite{karwowski2024goodhart,gururangan2018annotation} and binary success metrics under limited samples~\cite{snyder2026beyond}. Loss collapse in diffusion policies~\cite{dong2025conditioning} can produce near-zero actions; when these occur during safety evaluation, CAR is systematically inflated. Concurrent work such as HazardArena~\cite{chen2026hazardarena} expands the landscape to semantic safety. More broadly, any binary safety metric faces a base-rate problem when most episodes are collision-free through inaction~\cite{gururangan2018annotation}. dfCAR (\S\ref{sec:dfcar}) addresses this by conditioning on active behavior.

\paragraph{Positioning.}
The methods above share a common assumption: that explicit safety mechanisms---whether training-time losses, runtime filters, or constrained optimization---are necessary for collision avoidance in VLA policies. Our findings challenge this assumption: on SafeLIBERO, collision-free demonstrations alone produce avoidance equal to the best geometric supervision, suggesting that the bottleneck is data quality, not loss design.
